\documentclass[xcolor=dvipsnames,table]{beamer}

\usepackage{latexsym}
\usepackage[utf8]{inputenc}
\usepackage[brazil]{babel}
\usepackage{amssymb}
\usepackage{amsmath}
\usepackage{stmaryrd}
\usepackage{fancybox}
\usepackage{datetime}
\usepackage[T1]{fontenc}
\usepackage{graphicx}
\usepackage{graphics}
\usepackage{url}
\usepackage{algorithmic}
\usepackage{algorithm}
\usepackage{acronym}
\usepackage{array}

\usepackage{listings}
\usepackage{color}

\definecolor{mygreen}{rgb}{0,0.6,0}
\definecolor{mygray}{rgb}{0.5,0.5,0.5}
\definecolor{mymauve}{rgb}{0.58,0,0.82}

\lstdefinelanguage{JavaScript}{
  keywords={typeof, new, true, false, catch, function, return, null, catch, switch, var, if, in, while, do, else, case, break},
  keywordstyle=\color{blue}\bfseries,
  ndkeywords={class, export, boolean, throw, implements, import, this},
  ndkeywordstyle=\color{darkgray}\bfseries,
  identifierstyle=\color{black},
  sensitive=false,
  comment=[l]{//},
  morecomment=[s]{/*}{*/},
  commentstyle=\color{purple}\ttfamily,
  stringstyle=\color{red}\ttfamily,
  morestring=[b]',
  morestring=[b]",
}

\lstset{ %
  backgroundcolor=\color{white},   % choose the background color; you must add \usepackage{color} or \usepackage{xcolor}
  basicstyle=\small,        % the size of the fonts that are used for the code
  breakatwhitespace=false,         % sets if automatic breaks should only happen at whitespace
  breaklines=true,                 % sets automatic line breaking
  captionpos=b,                    % sets the caption-position to bottom
  commentstyle=\color{mygreen},    % comment style
  deletekeywords={...},            % if you want to delete keywords from the given language
  escapeinside={\%*}{*)},          % if you want to add LaTeX within your code
  extendedchars=true,              % lets you use non-ASCII characters; for 8-bits encodings only, does not work with UTF-8
  frame=single,	                   % adds a frame around the code
  keepspaces=true,                 % keeps spaces in text, useful for keeping indentation of code (possibly needs columns=flexible)
  keywordstyle=\color{blue},       % keyword style
  language=HTML,                 % the language of the code
  otherkeywords={*,...},           % if you want to add more keywords to the set
  numbers=left,                    % where to put the line-numbers; possible values are (none, left, right)
  numbersep=5pt,                   % how far the line-numbers are from the code
  numberstyle=\tiny\color{mygray}, % the style that is used for the line-numbers
  rulecolor=\color{black},         % if not set, the frame-color may be changed on line-breaks within not-black text (e.g. comments (green here))
  showspaces=false,                % show spaces everywhere adding particular underscores; it overrides 'showstringspaces'
  showstringspaces=false,          % underline spaces within strings only
  showtabs=false,                  % show tabs within strings adding particular underscores
  stepnumber=1,                    % the step between two line-numbers. If it's 1, each line will be numbered
  stringstyle=\color{mymauve},     % string literal style
  tabsize=2,	                   % sets default tabsize to 2 spaces
  title=\lstname,                   % show the filename of files included with \lstinputlisting; also try caption instead of title
  moredelim=**[is][\color{purple}]{@}{@},
}

\newtheorem{definicao}{Definio}
\newcommand{\tab}{\hspace*{2em}}

\mode<presentation>
{
  \definecolor{colortexto}{RGB}{0,0,0}
 
  \setbeamertemplate{background canvas}[vertical shading][ bottom=white!10,top=white!10]
  \setbeamercolor{normal text}{fg=colortexto} 

  \usetheme{Warsaw}
}

\title{HTML, CSS e JavaScript (Parte 2)} 

\author{
  Esdras Lins Bispo Jr. \\ \url{bispojr@ufg.br}
  } 
 \institute{
  Física para Ciência da Computação \\Bacharelado em Ciência da Computação}
\date{\textbf{06 de outubro de 2016} }

\logo{\includegraphics[width=1cm]{images/ufgJataiLogo.png}}

\begin{document}

	\begin{frame}
		\titlepage
	\end{frame}

	\AtBeginSection{
		\begin{frame}{Sumário}%[allowframebreaks]{Sumário}
    		\tableofcontents[currentsection]
    		%\tableofcontents[currentsection, hideothersubsections]
		\end{frame}
	}

	\begin{frame}{Plano de Aula}
		\tableofcontents
		%\tableofcontents[hideallsubsections]
	\end{frame}

%------------------------------------------
	\section{Revisão}
		
\subsection{HTML5}
\begin{frame}[fragile]{HTML5}
	\begin{block}{O que é?}
			HTML5 é a nova ``encarnação'' do padrão HTML. \\Ele traz novas capacidades para o navegador.
		\end{block} 
		\begin{block}{Um exemplo mínimo em HTML5}
			\begin{lstlisting}
<!doctype html>
<html>
	<head>
		<meta charset="utf-8">
		<title>Um documento minimo em HTML5</title>
	</head>
	<body>
		<h1>Boa noite, HTML5!</h1>
	</body>
</html>
\end{lstlisting}
		\end{block}
\end{frame}

\begin{frame}[fragile]{HTML5}
	\begin{block}{Elemento Canvas}
		Permite renderizar gráficos e animação no navegador sem necessitar de uso de plugins externos (como {\it FlashPlayer}).
	\end{block} 
		\begin{block}{...basta incluir esta linha no corpo do documento}
			\begin{lstlisting}
<canvas id="canvas" width="700" height="500">
</canvas>
\end{lstlisting}
		\end{block}
\end{frame}

\begin{frame}[fragile]{HTML5}
	\begin{block}{Exemplo}
		\begin{lstlisting}
<!doctype html>
<html>
	<head>
		<meta charset="utf-8">
		<title>Um documento minimo em HTML5</title>
	</head>
	<body>
		<canvas id="canvas" width="700" height="500">
		</canvas>
	</body>
</html>
\end{lstlisting}
		\end{block}
\end{frame}

\begin{frame}[fragile]{HTML5}
	\begin{block}{Pôr estilo no Canvas}
		Como qualquer elemento HTML. Você pode, por exemplo, ligar o arquivo HTML a um arquivo CSS externo.
	\end{block}
	\begin{block}{...basta incluir esta linha no {\tt head} do documento}
			\begin{lstlisting}
<link rel="stylesheet" href="style.css">
\end{lstlisting}
	\end{block}
\end{frame}

\begin{frame}[fragile]{HTML5}
	\begin{block}{Exemplo}
		\begin{lstlisting}
<!doctype html>
<html>
	<head>
		<meta charset="utf-8">
		<title>Um documento minimo em HTML5</title>
		<link rel="stylesheet" href="style.css">
	</head>
	<body>
		<canvas id="canvas" width="700" height="500">
		</canvas>
	</body>
</html>
\end{lstlisting}
		\end{block}
\end{frame}

\subsection{JavaScript}
\begin{frame}[fragile]{Adicionando JavaScript}
	\begin{block}{Ligando o arquivo .js ao corpo do HTML}
		\begin{lstlisting}
<!doctype html>
<html>
	<head>
		<meta charset="utf-8">
		<title>Um documento minimo em HTML5</title>
		<link rel="stylesheet" href="style.css">
	</head>
	<body>
		<canvas id="canvas" width="700" height="500">
		</canvas>
		<script src= "bolaQuicando.js"></script>
	</body>
</html>
\end{lstlisting}
		\end{block}
\end{frame}

\begin{frame}{Depurando no JavaScript}
	\begin{block}{Via Google Chrome...}
		Tecle \fbox{Ctrl + Shift + J}.
	\end{block} 
	\begin{block}{Teste...}
		\begin{itemize}
			\item {\tt 2 + 3}
			\item {\tt console.log("Hello World")}
			\item {\tt a=2; b=3; console.log(a*b);}
		\end{itemize}
	\end{block}
\end{frame}

\begin{frame}{Objetos em JavaScript}
	\begin{block}{Considerações iniciais...}
		Em JavaScript, os objetos são as unidades básicas.
	\end{block} 
	\begin{block}{O que é um objeto?}
		Uma coleção de propriedades	
	\end{block} 
	\begin{block}{O que é uma propriedade?}
		É uma associação entre um nome e um valor.
	\end{block}
	\begin{block}{O que é um valor?}
		Pode ser um número inteiro, uma cadeia, uma função, entre outros.
	\end{block}
\end{frame}

\begin{frame}{Objetos predefinidos}
	\begin{block}{Alguns Objetos predefinidos...}
		\begin{itemize}
			\item {\tt String}
			\item {\tt Array}
			\item {\tt Date}
			\item {\tt Math}...
		\end{itemize}
	\end{block} 
	\begin{block}{Criando novos objetos...}
		\begin{itemize}
			\item {\tt obj = new Object();}
			\\ou
			\item {\tt obj = \{\};}
		\end{itemize}
	\end{block}
\end{frame}

\begin{frame}[fragile]{Atribuindo propriedades}
	\begin{block}{Exemplo 1: Utilizando ponto...}
		\begin{lstlisting}[language=JavaScript]
obj = new Object();
obj.nome = "Primeiro objeto";
obj.length = 20;
console.log(obj.nome, obj.length);
\end{lstlisting}	
	\end{block} 
	\begin{block}{Exemplo 2: Utilizando colchetes...}
		\begin{lstlisting}[language=JavaScript]
obj = {};
obj["nome"] = "Primeiro objeto";
obj["length"] = 20;
\end{lstlisting}	
	\end{block}
\end{frame}

\begin{frame}[fragile]{Funções e Métodos}
	\begin{block}{O que é uma função?}
		É um bloco de código que é executado quando chamado pelo seu nome.
	\end{block} 
	\begin{block}{Sintaxe para a definição de função sem argumentos}
		\begin{lstlisting}[language=JavaScript]
function nomeDaFuncao(){
	//bloco de codigo
}
\end{lstlisting}	
	\end{block}
\end{frame}

\begin{frame}[fragile]{Funções e Métodos}
	\begin{block}{Sintaxe para a definição de função com argumentos}
		\begin{lstlisting}[language=JavaScript]
function nomeDaFuncao(arg1, arg2){
	//bloco de codigo
}
\end{lstlisting}	
	\end{block} 
	\begin{block}{A função pode retornar valores}
		\begin{lstlisting}[language=JavaScript]
function multiplicar(x, y){
	return x*y;
}
\end{lstlisting}	
	\end{block}
\end{frame}

\section{Continuação...}
\begin{frame}[fragile]{Funções e Métodos}
	\begin{block}{Atribuindo uma função como propriedade de um objeto...}
		\begin{lstlisting}[language=JavaScript]
nomeDoObjeto.nomeDaPropriedade = nomeDaFuncao;
\end{lstlisting}	
	\end{block} \pause
	\begin{block}{Exemplo}
		\begin{lstlisting}[language=JavaScript]
obj.multiplicar = multiplicar;
\end{lstlisting}	
	\end{block}
\end{frame}

\section{Herança, Construtores e Protótipos}
\subsection{Herança}
\begin{frame}{Herança}
	\begin{center}
		\includegraphics[scale=0.3]{images/heranca01.png}
	\end{center}
\end{frame}

\begin{frame}{Herança}
	\begin{center}
		\includegraphics[scale=0.3]{images/heranca02.png}
	\end{center}
\end{frame}

\begin{frame}{Herança}
	\begin{center}
		\includegraphics[scale=0.3]{images/heranca03.png}
	\end{center}
\end{frame}

\subsection{Construtores}
\begin{frame}[fragile]{Construtor}
	\begin{block}{Objeto a partir de uma função}
		É possível construir um objeto a partir de uma ``função-modelo''.
		\begin{lstlisting}[language=JavaScript]
function Particula(pnome){
	this.nome = pnome;
	this.mover = function(){
		console.log(this.nome + " em movimento!");
	};
}
\end{lstlisting}	
	\end{block}
\end{frame}

\begin{frame}[fragile]{Construtor}
	\begin{block}{Objeto a partir de uma função}
		\begin{lstlisting}[language=JavaScript]
particula1 = new Particula("eletron");
console.log(particula1.nome); //exibe "eletron"
particula1.mover();	
			//exibe "eletron em movimento"
\end{lstlisting}	
	\end{block} \pause
	\begin{block}{Construtor...}
		Se uma função é utilizada intencionalmente para criar novos objetos, a chamamos de {\bf construtor}.
	\end{block}
\end{frame}

\subsection{Protótipo}
\begin{frame}[fragile]{Protótipo}
	\begin{block}{Adicionando propriedades...}
		Após declarado, é possível adicionar propriedades a um construtor.
		\begin{lstlisting}[language=JavaScript]
Particula.prototype.massa = 1;
Particula.prototype.parar = 
		function(){console.log("Eu parei.");};
\end{lstlisting}	
	\end{block} \pause
	\begin{block}{Para novos objetos...}
		\begin{lstlisting}[language=JavaScript]
particula2 = new Particula("proton");
console.log(particula2.massa); //exibe 1
\end{lstlisting}	
	\end{block}
\end{frame}

\begin{frame}[fragile]{Protótipo}
	\begin{block}{Entretanto...}
		\begin{lstlisting}[language=JavaScript]
particula3 = new Particula("neutron");
particula3.massa = 2;	//propriedade com valor 2
console.log(Particula.prototype.massa);	//exibe 1
\end{lstlisting}	
	\end{block} \pause
	\begin{alertblock}{Necessário ter em mente...}
		A propriedade {\tt prototype} modifica o objeto-pai.	
	\end{alertblock}
\end{frame}

\begin{frame}[fragile]{Protótipo}
	\begin{block}{Algo interessante...}
		É possível adicionar propriedades apenas para o objeto-filho.
		\begin{lstlisting}[language=JavaScript]
particula4 = new Particula("neutrino");
particula4.spin = 0; //exibe 0
\end{lstlisting}	
	\end{block}
\end{frame}

\begin{frame}[fragile]{Exemplo...}
	\begin{block}{Objeto Bola}
		\begin{lstlisting}[language=JavaScript]
function Bola (raio, cor) {
	this.raio = raio;
	this.cor = cor;
	this.x = 0;
	this.y = 0;
	this.vx = 0;
	this.vy = 0;
}
\end{lstlisting}	
	\end{block}
\end{frame}

\begin{frame}[fragile]{Exemplo...}
	\begin{block}{Objeto Bola}
		\begin{lstlisting}[language=JavaScript]
Bola.prototype.desenhar = function (contexto) {
	contexto.fillStyle = this.cor;
	contexto.beginPath();
	contexto.arc(this.x, this.y, 
							this.raio, 0, 2*Math.PI, true);
	contexto.closePath();
	contexto.fill();
};
\end{lstlisting}	
	\end{block}
\end{frame}

\begin{frame}[fragile]{Exemplo...} 
	\begin{block}{HTML Canvas {\tt arc}}
		A propriedade {\tt arc()} cria um arco \\(utilizado para criar círculos ou parte de círculos). 
	\end{block}
	\begin{block}{Sintaxe}
		{\tt arc(x, y, raio, angIni, angFin, antiHor);}
	\end{block}
	\begin{block}{Argumentos}
		\begin{itemize}
			\item {\tt x} e {\tt y}: coordenadas do arco;
			\item {\tt raio}: raio do arco; 
			\item {\tt angIni} e {\tt angFin}: ângulos inicial e final do arco \\(em radianos);
			\item {\tt antiHor}: valor booleano para o sentido anti-horário.
		\end{itemize}
	\end{block}
\end{frame}

\begin{frame}{Exemplo...}
	\begin{center}
		\includegraphics[scale=0.4]{images/arc.png}
	\end{center}
\end{frame}

\begin{frame}[fragile]{Exemplo...}
	\begin{block}{Objeto Bola}
		\begin{lstlisting}[language=JavaScript]
var canvas = document.getElementById('canvas');
var contexto = canvas.getContext('2d');
var bola = new Bola(50,'#0000ff');
bola.x = 100;
bola.y = 100;
bola.desenhar(contexto);
\end{lstlisting}	
	\end{block}
\end{frame}

	
	\begin{frame}
		\titlepage
	\end{frame}
	
\end{document}